\section*{Chapter \ref{ch:g12:processes-and-scheduling} --- Processes and Scheduling}

\begin{solution}{exo:g12:processes-and-scheduling:1}
In order of arrival: X from $0$ to $4$, Y from $4$ to $5$, Z from $5$ to $8$. The Gantt line is \code{XXXXYZZZ}.
\end{solution}

\begin{solution}{exo:g12:processes-and-scheduling:2}
X waits $0$, Y waits $4$, Z waits $5$: a mean of $9/3 = 3$.
\end{solution}

\begin{solution}{exo:g12:processes-and-scheduling:3}
\code{XYZXZXZX}. Y is done after its first unit; Z after its third turn, at time $7$; X runs alone at the end.
\end{solution}

\begin{solution}{exo:g12:processes-and-scheduling:4}
A switch happens each time the next slice belongs to another \omterm{def:g12:processes-and-scheduling:process}{process}. With $q = 2$ the $10$ slices are A, B, C, A, D,
B, C, A, C, C: $8$ switches, since C's last two slices follow each other. With $q = 4$ the $6$ slices A, B, C, D, A, C
make $5$.
\end{solution}

\begin{solution}{exo:g12:processes-and-scheduling:5}
\smallskip\noindent\hspace*{2em}\code{def turnaround(workload, slices):}

\noindent\hspace*{4em}\code{finish = {name: end for name, _, end in slices}}

\noindent\hspace*{4em}\code{return {name: finish[name] - arrival for name, arrival, _ in workload}}

\smallskip\noindent The dictionary keeps the last slice of each name, which is where the \omterm{def:g12:processes-and-scheduling:process}{process} finishes.
\end{solution}

\begin{solution}{exo:g12:processes-and-scheduling:6}
Each \omterm{def:g12:processes-and-scheduling:process}{process} runs once, for at most $2$ units: A from $0$ to $2$, B from $2$ to $4$, C from $4$ to $6$, D from $6$ to
$8$. Then the queue is empty and the simulation stops. Only D, whose burst is $2$, has finished; A, B and C are left
with $3$, $1$ and $6$ units, and will never run again. The tests check these remainders.
\end{solution}

\begin{solution}{exo:g12:processes-and-scheduling:7}
$20$. An interleaving of $6$ steps is fixed by choosing which $3$ of the $6$ places hold P's steps, in their own
order; the other places hold Q's. There are $\binom{6}{3} = 20$ such choices. For $a$ and $b$ steps:
$\binom{a+b}{a}$, which is how \code{interleavings} enumerates them, with \code{combinations}.
\end{solution}

\begin{solution}{exo:g12:processes-and-scheduling:8}
\smallskip\noindent\hspace*{2em}\code{def sjf(workload):}

\noindent\hspace*{4em}\code{pending, time, slices = sorted(workload, key=lambda w: w[1]), 0, []}

\noindent\hspace*{4em}\code{while pending:}

\noindent\hspace*{6em}\code{arrived = [w for w in pending if w[1] <= time]}

\noindent\hspace*{6em}\code{if not arrived:}

\noindent\hspace*{8em}\code{time = pending[0][1]}

\noindent\hspace*{8em}\code{continue}

\noindent\hspace*{6em}\code{job = min(arrived, key=lambda w: (w[2], w[1]))}

\noindent\hspace*{6em}\code{pending.remove(job)}

\noindent\hspace*{6em}\code{slices.append((job[0], time, time + job[2]))}

\noindent\hspace*{6em}\code{time += job[2]}

\noindent\hspace*{4em}\code{return slices}

\smallskip\noindent On the chapter's workload, A is alone at time $0$ and runs to $5$; then D ($2$), B ($3$) and C
($8$): D from $5$ to $7$, B to $10$, C to $18$. Turnarounds $5$, $9$, $16$, $4$: a mean of $8.5$, against $10.25$ for
first come, first served.
\end{solution}

\begin{solution}{exo:g12:processes-and-scheduling:9}
The \omterm{def:g12:processes-and-scheduling:race}{lock} stays held by a \omterm{def:g12:processes-and-scheduling:process}{process} that has left the \omterm{def:g12:processes-and-scheduling:race}{critical section}. Every other \omterm{def:g12:processes-and-scheduling:process}{process} that tries to acquire it
waits for ever, and so does the same \omterm{def:g12:processes-and-scheduling:process}{process} on its next call. The fix releases the \omterm{def:g12:processes-and-scheduling:race}{lock} on every path: in Python,
\code{with lock:} around the \omterm{def:g12:processes-and-scheduling:race}{critical section}, which releases it however the block is left, even by an exception.
\end{solution}

\begin{solution}{exo:g12:processes-and-scheduling:10}
P loads $0$; Q loads $0$, adds, stores $1$; P adds to its own copy, $1$, and stores $1$. The final counter is $1$: Q's
update is lost, overwritten by P's.
\end{solution}

\begin{solution}{exo:g12:processes-and-scheduling:11}
The first design rules out \omterm{def:g12:processes-and-scheduling:race}{mutual exclusion} as far as the processes are concerned: none of them ever holds the
printer, only the server does, so no \omterm{def:g12:processes-and-scheduling:process}{process} can wait for another to give it back. The second rules out hold and
wait: a \omterm{def:g12:processes-and-scheduling:process}{process} that waits holds nothing, and a \omterm{def:g12:processes-and-scheduling:process}{process} that holds its resources never waits for more.
\end{solution}

\begin{solution}{exo:g12:processes-and-scheduling:12}
Between its \code{acquire} and its \code{release}, a \omterm{def:g12:processes-and-scheduling:process}{process} holds the \omterm{def:g12:processes-and-scheduling:race}{lock}, and the other cannot pass its own
\code{acquire}. So the three steps of one \omterm{def:g12:processes-and-scheduling:process}{process} are never interleaved with those of the other: one \omterm{def:g12:processes-and-scheduling:process}{process}, say the
first to acquire, does load, add, store, and releases; only then can the second load. Two cases, P first or Q first,
and in both the second \omterm{def:g12:processes-and-scheduling:process}{process} loads the value $1$ stored by the first and stores $2$. Steps outside the \omterm{def:g12:processes-and-scheduling:race}{lock} are
only \code{acquire} and \code{release}, which do not touch the counter. The program finds the same: exactly $2$ runs,
both ending at $2$.
\end{solution}

\begin{solution}{exo:g12:processes-and-scheduling:13}
A depth-first search that marks the vertices still being explored: meeting one of them again closes a cycle.

\smallskip\noindent\hspace*{2em}\code{def has_cycle(waits_for):}

\noindent\hspace*{4em}\code{state = {}}

\noindent\hspace*{4em}\code{def visit(v):}

\noindent\hspace*{6em}\code{state[v] = "open"}

\noindent\hspace*{6em}\code{for w in waits_for.get(v, []):}

\noindent\hspace*{8em}\code{if state.get(w) == "open" or (w not in state and visit(w)):}

\noindent\hspace*{10em}\code{return True}

\noindent\hspace*{6em}\code{state[v] = "done"}

\noindent\hspace*{6em}\code{return False}

\noindent\hspace*{4em}\code{return any(v not in state and visit(v) for v in list(waits_for))}

\smallskip\noindent A vertex marked \code{done} has no cycle below it, so it is never explored twice. The tests try a
cycle of three processes and a chain without one.
\end{solution}

\begin{solution}{exo:g12:processes-and-scheduling:14}
First come, first served: waits $0$, $6$, $14$, $17$, a mean of $37/4 = 9.25$. Shortest job first runs $2$, $3$, $6$,
$8$: waits $0$, $2$, $5$, $11$, a mean of $18/4 = 4.5$.

When all arrive at $0$, the \omterm{def:g12:processes-and-scheduling:process}{process} in place $k$ (from $1$ to $n$) makes the $n - k$ processes after it wait for its
burst, so the total waiting time is the sum of each burst times the number of processes after it. Take any order in
which a longer burst $b$ comes just before a shorter one $c < b$. Swapping the two changes nothing for the others,
and changes the total by $c - b < 0$: the swap improves the order. So an order with the smallest total has no such
pair: its bursts are increasing, which is shortest job first. The tests compare it with all $24$ orders.
\end{solution}

\begin{solution}{exo:g12:processes-and-scheduling:15}
Number the philosophers and forks $0$ to $4$; philosopher $k$ uses forks $k$ and $k + 1$ (fork $0$ for the last).
If all five take their left fork at once, every fork is held and each waits for the right one, held by a neighbour:
a cycle of five, a \omterm{def:g12:processes-and-scheduling:deadlock}{deadlock}. With the rule ``lower-numbered fork first'', philosophers $0$ to $3$ still start on the
left, but philosopher $4$ starts with fork $0$. A cycle would need each philosopher to hold one fork and wait for a
higher-numbered one; philosopher $4$ would hold fork $0$ and wait for fork $4$, but then philosopher $0$ holds
nothing, and the cycle breaks. The function \code{philosophers} tries every reachable state and confirms: a \omterm{def:g12:processes-and-scheduling:deadlock}{deadlock}
without the rule, none with it.
\end{solution}

\begin{solution}{pb:g12:processes-and-scheduling:1}
\textbf{1.} $20 + 3 + 6 + 2 + 4 = 35$ units.

\textbf{2.} The core does one unit of work per unit of time, and the work starts at $0$, so at least $35$; five
processes also need at least four switches, so at least $39$.

\textbf{3.} By arrival, with a switch between each: A from $0$ to $20$, B $21$--$24$, C $25$--$31$, D $32$--$34$,
E $35$--$39$.

\textbf{4.} $20$, $23$, $29$, $31$, $34$: $137 / 5 = 27.4$, the value printed.

\textbf{5.} Each unit of work is followed by a switch, so $22$ units are lost. B, D and E finish sooner than under
first come, first served (turnarounds $20$, $16$, $30$ against $23$, $31$, $34$), but C, sharing the core one unit at
a time, takes $41$ instead of $29$, and A $57$ instead of $20$. The total grows by $27$ units, the mean by $5.4$.

\textbf{6.} $22$ switches of $1$ unit, and $35 + 22 = 57$, the time printed: the core is never idle.

\textbf{7.} With $q \ge 20$, every \omterm{def:g12:processes-and-scheduling:process}{process} finishes in its first turn, so nobody returns to the queue: the processes
run once each, in order of arrival, which is first come, first served. The tests check $q = 20$ to $29$.

\textbf{8.} With $q = 12$, A is interrupted at $12$ and goes behind B, C, D and E, who all finish in their first turn;
A comes back for its last $8$ units. The short processes finish $8$ units sooner; the extra switch makes everything
end at $40$.

\textbf{9.} $q = 6$, with $20.2$: $7.2$ less than first come, first served.

\textbf{10.} A $0$--$6$, B $7$--$10$, C $11$--$17$, D $18$--$20$, E $21$--$25$, then A alone from $26$ to $40$. B, C, D
and E finish at $10$, $17$, $20$, $25$; A at $40$.

\textbf{11.} $6$ is the longest burst of the short processes ($3$, $6$, $2$, $4$): each of them finishes in one turn,
while A, the long one, is interrupted once and goes to the back.

\textbf{12.} Every short \omterm{def:g12:processes-and-scheduling:process}{process} already finishes in one turn, so a longer \omterm{def:g12:processes-and-scheduling:scheduler}{quantum} only lengthens A's first slice:
B, C, D and E each finish one unit later per unit of \omterm{def:g12:processes-and-scheduling:scheduler}{quantum}. The total grows by $4$, the mean by $0.8$, as printed.

\textbf{13.} Finishing times for $q = 4$: A $42$, B $8$, C $29$, D $16$, E $26$; for $q = 5$: A $42$, B $9$, C $31$,
D $18$, E $23$. A's first slice is one unit longer, which delays B by $1$; C needs a second turn in both cases, and
everything before it is longer, so C and D lose $2$. E gains $3$: it arrives at $5$, which with $q = 5$ is exactly
when A's slice ends, so E joins the queue before A; with $q = 4$, A is back in the queue at $4$, ahead of E. In all
$+1 + 2 + 2 - 3 = +2$ units, so the mean grows by $0.4$.

\textbf{14.} First come, first served: B arrives at $1$ and starts at $21$, a response time of $20$. With $q = 6$:
it starts at $7$, a response of $6$.

\textbf{15.} The typist wants short response times, so a short \omterm{def:g12:processes-and-scheduling:scheduler}{quantum}: with $q = 1$ every \omterm{def:g12:processes-and-scheduling:process}{process} starts within a
few units. The person waiting for a report wants the short turnarounds of $q = 6$.

\textbf{16.} Without switches, the $22$ lost units disappear, and \omterm{def:g12:processes-and-scheduling:scheduler}{round robin} with $q = 1$ shares the core finely,
which lets the short processes finish early. It is still not best: the four short ones are slowed down by sharing
with each other, a little more than with $q = 2$ or $3$, where each takes longer turns.

\textbf{17.} No: with the same switch cost of $1$, shortest job first gives $26.0$ ($24.0$ if switches were free). A is alone at time $0$ and, once started, runs to its end; the
short processes all wait $20$ units. \omterm{def:g12:processes-and-scheduling:scheduler}{Round robin} \omterm{def:g12:processes-and-scheduling:interrupt}{interrupts} A, which shortest job first never does. A real \omterm{def:g12:processes-and-scheduling:scheduler}{scheduler}
cannot apply it as written because it does not know the bursts in advance: they depend on what the programs do.

\textbf{18.} $4 / 4004 \approx 0.1\,\%$ of the time. Here, with $q = 6$ and a switch of $1$, up to $1/7$, about
$14\,\%$: the toy cost is more than a hundred times heavier.

\textbf{19.} A \omterm{def:g12:processes-and-scheduling:process}{process} that blocks often, waiting for a key or the network, uses little of its \omterm{def:g12:processes-and-scheduling:scheduler}{quantum}, and is ready
again as soon as its event comes. Serving it first keeps its response short, and costs the long computations little,
since it soon gives the core back.

\textbf{20.} \textbf{For five processes, one of $20$ units and four of $2$ to $6$, with a switch costing one unit,
\omterm{def:g12:processes-and-scheduling:scheduler}{round robin} with a \omterm{def:g12:processes-and-scheduling:scheduler}{quantum} of $6$ gives the smallest mean turnaround, $20.2$ units, against $27.4$ for first come,
first served and $32.8$ for a \omterm{def:g12:processes-and-scheduling:scheduler}{quantum} of $1$.}
\end{solution}
